\chapter{Experimentos y resultados}
\label{cap:resultados}

Este capítulo presenta el diseño experimental y los resultados obtenidos.
\textit{(Este capítulo debe completarse con los resultados reales a medida que
se ejecuten los experimentos; a continuación se proporciona la estructura y las
plantillas de tablas y figuras.)}

\section{Diseño experimental}
Se definen los siguientes experimentos, alineados con los objetivos O4 y O5:
\begin{description}
  \item[E1.] Clasificación con descriptores topológicos + clasificador clásico.
  \item[E2.] Clasificación con CNN (\textit{transfer learning}).
  \item[E3.] Fusión de características topológicas y \textit{embeddings} de CNN.
  \item[E4.] Análisis de robustez frente a ruido y a variaciones de adquisición.
\end{description}

\section{Configuración}
Se detallan los hiperparámetros, la partición de datos por paciente, el número
de repeticiones y el hardware utilizado.

\section{Resultados}

\begin{table}[H]
\centering
\caption{Resultados comparativos en el conjunto de test (plantilla).}
\label{tab:resultados}
\begin{tabular}{@{}lccccc@{}}
\toprule
\textbf{Modelo} & \textbf{Exactitud} & \textbf{Sensib.} & \textbf{Especif.} &
\textbf{F1} & \textbf{AUC} \\
\midrule
Baseline (clase mayoritaria) & -- & -- & -- & -- & -- \\
TDA + SVM                    & -- & -- & -- & -- & -- \\
TDA + Random Forest          & -- & -- & -- & -- & -- \\
CNN (transfer learning)      & -- & -- & -- & -- & -- \\
Fusión TDA + CNN             & -- & -- & -- & -- & -- \\
\bottomrule
\end{tabular}
\end{table}

\begin{figure}[H]
\centering
% \includegraphics[width=0.7\textwidth]{figuras/roc.png}
\fbox{\parbox[c][4cm][c]{0.7\textwidth}{\centering\itshape
Curvas ROC de los distintos modelos. (Sustituir por la figura definitiva.)}}
\caption{Comparación de curvas ROC.}
\label{fig:roc}
\end{figure}

\section{Discusión}
Se analizan los resultados a la luz de los objetivos: capacidad discriminativa de
los descriptores topológicos, comparación con el \textit{deep learning},
beneficio de la fusión, robustez frente al ruido e interpretabilidad. Se discuten
también las limitaciones observadas.
